%! Unit 1: Asking Questions with Data: Study Design — Summative Assessment — Practice Test --- Study Copy (blank)
% Rewritten 2026-09-28 against the regenerated lessons 1.0-1.8: every item now aims at the
% target misconception its lesson names. Parallel to ../actual_test/ item for item.
%
% THE KEY IS GENERATED, NEVER HAND-EDITED:
%   python3 templates/lesson/mkkey.py unit01/tests/practice_test/main.tex \
%       unit01/test_keys/practice_test_key/main.tex unit01/test_keys/practice_test_key/answers.json
% Re-run it after every edit to this file.
\documentclass[10pt]{article}
\usepackage{saar-article}
\usepackage{saar-boxes}

% Work blocks on a test need handwriting room, not typeset spacing. This moves
% the blank and the key together, so it cannot open a page mismatch.
\setlength{\workrowsep}{5pt}

% Part-divider strip
\newcommand{\parthead}[1]{%
  \vspace{6pt}\noindent
  \begin{headlinebox}{cerulean}{\color{white}\bfseries #1}\end{headlinebox}%
  \vspace{4pt}}

% The correct multiple-choice option. Prints nothing in the blank; prints the marker once
% -key is loaded (the same switch that makes `work` blocks visible), so the key regenerates
% from this file byte for byte.
\makeatletter
\newcommand{\correct}{\ifsaar@workvisible\ \textcolor{keyred}{\textbf{$\leftarrow$ correct}}\fi}
% Key-only drawing (the correct bars on item 30), inside the figure's existing bounding box.
\newcommand{\keyonly}[1]{\ifsaar@workvisible#1\fi}
\makeatother

\begin{document}

\pageheader{Unit 1: Asking Questions with Data: Study Design}{Practice Test --- Study Copy}
\namedateperiod

\vspace{0.10in}
\begin{remindbox}
\small
\textbf{This is a practice test.} It mirrors the real Unit 1 test in length, format,
and the ideas it covers, but the numbers and contexts differ. Work every problem
with no notes, then check your answers against the key.
\textbf{100 points:} Part A $14$, Part B $16$, Part C $40$, Part D $30$.
\end{remindbox}

% ======================================================================
\parthead{Part A --- Vocabulary \hfill 14 questions, 1 point each}

\small
\textbf{Set 1 --- Data and Sampling.} Write the letter of the matching description
next to each term.

\vspace{2pt}
\begin{minipage}[t]{0.36\linewidth}
\begin{enumerate}[leftmargin=1.9em, itemsep=3.5pt, topsep=1pt]
  \item \blank{0.8cm} Categorical variable
  \item \blank{0.8cm} Statistical question
  \item \blank{0.8cm} Parameter
  \item \blank{0.8cm} Statistic
  \item \blank{0.8cm} Sampling frame
  \item \blank{0.8cm} Stratified sample
  \item \blank{0.8cm} Cluster sample
\end{enumerate}
\end{minipage}\hfill
\begin{minipage}[t]{0.61\linewidth}
\begin{enumerate}[label=\Alph*., leftmargin=1.9em, itemsep=3.5pt, topsep=1pt]
  \item A number computed from a sample.
  \item The population is split into groups that are alike inside, and some are
        chosen at random from every group.
  \item A question whose answers are expected to vary, and that names the group
        and what will be recorded.
  \item A few whole groups, each a small mix of the whole, are chosen at random
        and everyone in them is studied.
  \item A variable whose values are labels or group names.
  \item The list of individuals a study can actually choose from.
  \item A number that describes a population.
\end{enumerate}
\end{minipage}

\vspace{6pt}
\textbf{Set 2 --- Bias and Study Design.} Same directions.

\vspace{2pt}
\begin{minipage}[t]{0.36\linewidth}
\begin{enumerate}[leftmargin=1.9em, itemsep=3.5pt, topsep=1pt, start=8]
  \item \blank{0.8cm} Undercoverage
  \item \blank{0.8cm} Nonresponse
  \item \blank{0.8cm} Response bias
  \item \blank{0.8cm} Lurking variable
  \item \blank{0.8cm} Confounding variable
  \item \blank{0.8cm} Placebo
  \item \blank{0.8cm} Randomized block design
\end{enumerate}
\end{minipage}\hfill
\begin{minipage}[t]{0.61\linewidth}
\begin{enumerate}[label=\Alph*., leftmargin=1.9em, itemsep=3.5pt, topsep=1pt]
  \item Who asks, the wording, or the order of the questions pushes people away
        from their true answer.
  \item A variable that differs between an experiment's groups and affects the
        response, so two explanations cannot be told apart.
  \item Units are sorted into groups of similar units first, then randomly
        assigned within each group.
  \item Part of the population is left off the frame and can never be chosen.
  \item A fake treatment built to look real.
  \item A variable outside the study that affects both variables being compared.
  \item Individuals were chosen for the sample but never answered.
\end{enumerate}
\end{minipage}
\normalsize

% ======================================================================
\boxguard[8]
\parthead{Part B --- Multiple Choice \hfill 8 questions, 2 points each}

\small
\begin{enumerate}[leftmargin=2.2em, itemsep=5pt, topsep=2pt, start=15]

\item Granite Falls Middle School records each student's \textbf{bus route number}
($4$, $7$, $12$, or $15$). This variable is
\begin{enumerate}[label=(\Alph*), leftmargin=1.8em, itemsep=1pt, topsep=1pt]
  \item quantitative, because the values are numbers.
  \item quantitative, because you can compute a mean route number.
  \item categorical, because the numbers are labels for routes.\correct
  \item neither, because a route number is not really data.
\end{enumerate}

% keep each stem with its four options
\boxguard[7]
\item Two volunteers each take an honest random sample of $50$ Silver Lake
Recreation Center members. One finds that $34\%$ swim laps weekly; the other finds
$42\%$. The best explanation is that
\begin{enumerate}[label=(\Alph*), leftmargin=1.8em, itemsep=1pt, topsep=1pt]
  \item both are statistics estimating one parameter, and samples vary.\correct
  \item one volunteer made a mistake, because there is only one true percent.
  \item the true percent must be exactly $38\%$, halfway between them.
  \item the percent of members who swim changed between the two samples.
\end{enumerate}

% keep each stem with its four options
\boxguard[7]
\item Each of Maplewood High School's $32$ homerooms holds students from a
\textbf{single grade}. A researcher draws $2$ homerooms at random and surveys
everyone in them. The main problem with this sample is that
\begin{enumerate}[label=(\Alph*), leftmargin=1.8em, itemsep=1pt, topsep=1pt]
  \item no chance was used, so it is a convenience sample.
  \item two homerooms is too small a sample to use.
  \item a cluster sample can never be used in a school.
  \item a homeroom is alike inside, so only one or two grades are reached.\correct
\end{enumerate}

% keep each stem with its four options
\boxguard[7]
\item Brookside Apartments draws $60$ households at random from its complete list
and mails each a survey. $22$ are never returned. The missing $22$ are an example of
\begin{enumerate}[label=(\Alph*), leftmargin=1.8em, itemsep=1pt, topsep=1pt]
  \item undercoverage, because those households were left out.
  \item nonresponse, because they were chosen but never answered.\correct
  \item a convenience sample, because the others were easier to reach.
  \item sampling variability, because samples differ from each other.
\end{enumerate}

% keep each stem with its four options
\boxguard[7]
\item A survey's method causes it to \textbf{underestimate} by $12$ percentage
points. The researcher repeats it with \textbf{five times as many people} and the
same method. The new estimate will most likely
\begin{enumerate}[label=(\Alph*), leftmargin=1.8em, itemsep=1pt, topsep=1pt]
  \item underestimate by about the same $12$ points.\correct
  \item land close to the true value, because the sample is larger.
  \item overestimate, because larger samples pull the other way.
  \item be impossible to predict in any direction.
\end{enumerate}

% keep each stem with its four options
\boxguard[7]
\item In a random sample of adults, those who own a dog walk more minutes per day
than those who do not. Nobody was assigned a dog. This result shows
\begin{enumerate}[label=(\Alph*), leftmargin=1.8em, itemsep=1pt, topsep=1pt]
  \item that owning a dog causes people to walk more.
  \item that walking more causes people to buy a dog.
  \item an association only; a lurking variable could explain it.\correct
  \item nothing at all, because the sample was chosen at random.
\end{enumerate}

% keep each stem with its four options
\boxguard[7]
\item In an experiment on a new vitamin, the control group is given a
\textbf{placebo}. Its purpose is to
\begin{enumerate}[label=(\Alph*), leftmargin=1.8em, itemsep=1pt, topsep=1pt]
  \item make the total number of subjects larger.
  \item let both groups believe they are treated, so only the vitamin differs.\correct
  \item tell each subject which group they are in.
  \item remove the need for random assignment.
\end{enumerate}

% keep each stem with its four options
\boxguard[7]
\item A researcher wants to know whether energy drinks cause heart problems in
$12$-year-olds. No experiment is run. The best reason is that the experiment would be
\begin{enumerate}[label=(\Alph*), leftmargin=1.8em, itemsep=1pt, topsep=1pt]
  \item impractical, because $12$-year-olds are hard to find.
  \item impossible, because heart problems cannot be measured.
  \item unnecessary, because observational studies prove causation.
  \item unethical, because it would assign children to drink them.\correct
\end{enumerate}

\end{enumerate}
\normalsize

% ======================================================================
\boxguard[8]
\parthead{Part C --- Short Answer \& Computation \hfill 8 questions, 5 points each}

\small
\begin{enumerate}[leftmargin=2.2em, itemsep=9pt, topsep=2pt, start=23]

% ---------------------------------------------------------------- C1 (1.0/1.1)
\item \textbf{Granite Falls Middle School} has $840$ students. For each student the
office records: bus route number, hours of sleep last night, favorite elective, and
number of siblings.

\vspace{2pt}
(a) Who are the \textbf{individuals} in this data set? \blank{6.2cm}

\vspace{3pt}
\boxguard[10]
(b) Classify each variable.

\vspace{2pt}
\renewcommand{\arraystretch}{1.25}
\begin{tabularx}{\linewidth}{@{} X c @{}}
\rowcolor{frost}
\textbf{Variable} & \textbf{Categorical or quantitative?} \\
\midrule
Bus route number & \blank{4.4cm} \\
Hours of sleep last night & \blank{4.4cm} \\
Favorite elective & \blank{4.4cm} \\
Number of siblings & \blank{4.4cm} \\
\end{tabularx}

\vspace{4pt}
(c) The office's first draft question is \emph{``How much sleep?''} Explain why that
is not a statistical question, then rewrite it so that it is.

\par\writeline
\par\writeline

% ---------------------------------------------------------------- C2 (1.2)
\boxguard[16]
\item \textbf{Silver Lake Recreation Center} has $1{,}600$ members. The director
surveys $80$ members chosen at random; $28$ of them use the pool every week.

\vspace{2pt}
(a) Population: \blank{4.4cm} \quad Sample size: \blank{1.8cm}

\vspace{3pt}
Sample: \blank{9.0cm}

\vspace{3pt}
(b) What percent of the sample uses the pool weekly? About how many of all $1{,}600$
members does that estimate?

\begin{work}
  \dfrac{28}{80} &= 0.35 = 35\% \\
  0.35 \times 1600 &= 560 \text{ members}
\end{work}

(c) Is the $35\%$ a parameter or a statistic? Say how you know.

\par\writeline

(d) Name one \textbf{constraint} that keeps the director from taking a census of all
$1{,}600$ members.

\par\writeline

% ---------------------------------------------------------------- C3 (1.3)
\boxguard[18]
\item \textbf{Maplewood High School} has $960$ students and wants a
\textbf{stratified} random sample of $64$, allocated in proportion to grade size.

\vspace{2pt}
(a) What fraction of the school is being sampled?

\begin{work}
  \dfrac{64}{960} &= \dfrac{1}{15}
\end{work}

(b) Show the work for the \textbf{10th grade} row, then complete the table.

\begin{work}
  255 \times \dfrac{1}{15} &= 17 \text{ students}
\end{work}

\vspace{-2pt}
\renewcommand{\arraystretch}{1.2}
\begin{tabularx}{\linewidth}{@{} X c c c c | c @{}}
\rowcolor{frost}
\textbf{Grade} & 9 & 10 & 11 & 12 & \textbf{Total} \\
\midrule
Students in grade & $300$ & $255$ & $225$ & $180$ & $960$ \\
Number sampled & \blank{1.1cm} & \blank{1.1cm} & \blank{1.1cm} & \blank{1.1cm} & \blank{1.1cm} \\
\end{tabularx}

% ---------------------------------------------------------------- C4 (1.3)
\boxguard[18]
\item \textbf{Brookside Apartments} has $720$ households on a numbered list.
The manager wants a \textbf{systematic} random sample of $48$.

\vspace{2pt}
(a) Find the interval $k$.

\begin{work}
  k &= 720 \div 48 \\
  k &= 15
\end{work}

(b) The random start is household $7$. List the first four households selected:

\vspace{2pt}
\blank{1.6cm} \quad \blank{1.6cm} \quad \blank{1.6cm} \quad \blank{1.6cm}

\vspace{4pt}
(c) Name the sampling technique each manager used.

\vspace{2pt}
\renewcommand{\arraystretch}{1.25}
\begin{tabularx}{\linewidth}{@{} X c @{}}
\rowcolor{frost}
\textbf{What the manager did} & \textbf{Technique} \\
\midrule
Split the households by building age and drew some at random from every
building, in proportion. & \blank{3.4cm} \\
Drew $4$ of the $24$ floors at random, then surveyed every household on
those $4$. & \blank{3.4cm} \\
Stood in the lobby on Saturday morning and asked whoever walked past. & \blank{3.4cm} \\
\end{tabularx}

% ---------------------------------------------------------------- C5 (1.4)
\boxguard[20]
\item \textbf{Elmhurst Public Library} has $1{,}500$ card holders and wants to know
what fraction want Sunday hours. Two random phone samples were taken.

\vspace{1pt}
\textbf{Method 1.} $100$ card holders were asked, \emph{``Most libraries now open on
Sundays. Wouldn't you like Sunday hours too?''} $74$ said yes. \\
\textbf{Method 2.} $120$ card holders were asked, \emph{``Should the library add
Sunday hours?''} $54$ said yes.

\vspace{2pt}
(a) Find the percent saying yes for each method.

\begin{work}
  \dfrac{74}{100} &= 0.74 = 74\% \\
  \dfrac{54}{120} &= 0.45 = 45\%
\end{work}

(b) Estimate the number of all $1{,}500$ card holders each method points to.

\begin{work}
  0.74 \times 1500 &= 1110 \text{ card holders} \\
  0.45 \times 1500 &= 675 \text{ card holders}
\end{work}

(c) Both samples were random, yet one method is biased. Name it, name the
\textbf{type} of bias, and say which \textbf{direction} it pushes the estimate.

\par\writeline
\par\writeline

% ---------------------------------------------------------------- C6 (1.5)
\boxguard[20]
\item A random sample of $200$ Maplewood students is taken. Of the $80$ who play a
\textbf{school sport}, $60$ made the honor roll. Of the $120$ who do not, $60$ made
the honor roll. Nobody was assigned --- students chose.

\vspace{2pt}
(a) Find each group's honor-roll percent, and the gap between them.

\begin{work}
  \dfrac{60}{80} &= 0.75 = 75\% \\
  \dfrac{60}{120} &= 0.50 = 50\% \\
  75\% - 50\% &= 25 \text{ percentage points}
\end{work}

(b) Explanatory variable: \blank{4.6cm} \quad Response: \blank{4.6cm}

\vspace{3pt}
(c) Using all $200$ students, find the overall honor-roll percent and estimate how
many of all $960$ students that is.

\begin{work}
  \dfrac{60 + 60}{200} &= \dfrac{120}{200} = 0.60 = 60\% \\
  0.60 \times 960 &= 576 \text{ students}
\end{work}

(d) Maplewood requires passing grades to stay on a team. Name the \textbf{reversed
arrow} or a \textbf{lurking variable}, and say which way it pushes the numbers.

\par\writeline
\par\writeline

% ---------------------------------------------------------------- C7 (1.6)
\boxguard[20]
\item Maplewood now runs an \textbf{experiment}. Of $80$ volunteers, a random
number generator sends $40$ to a review session before the unit test and $40$ to
an ordinary study hall. Afterwards $30$ of the review group and $22$ of the study
hall group passed.

\vspace{2pt}
(a) Find each group's pass percent and the gap.

\begin{work}
  \dfrac{30}{40} &= 0.75 = 75\% \\
  \dfrac{22}{40} &= 0.55 = 55\% \\
  75\% - 55\% &= 20 \text{ percentage points}
\end{work}

(b) Treatment: \blank{4.6cm} \quad Control group: \blank{4.6cm}

\vspace{3pt}
(c) Instead, Maplewood \textbf{blocks} on the last unit test: $48$ volunteers
passed it and $32$ did not. Randomly split each block in half. How many go to each
group?

\begin{work}
  48 \div 2 &= 24 \text{ per group (passed block)} \\
  32 \div 2 &= 16 \text{ per group (did-not-pass block)}
\end{work}

% ---------------------------------------------------------------- C8 (1.8)
\boxguard[30]
\item The Maplewood Student Council asked a stratified sample of $80$ students how
to spend a grant. Complete the table, then complete the bar graph.

\vspace{2pt}
\renewcommand{\arraystretch}{1.25}
\begin{tabularx}{\linewidth}{@{} X c c @{}}
\rowcolor{frost}
\textbf{Choice} & \textbf{Count} & \textbf{Relative frequency} \\
\midrule
New library books     & $28$ & \blank{2.6cm} \\
Picnic tables         & $20$ & \blank{2.6cm} \\
Water bottle stations & $24$ & \blank{2.6cm} \\
Longer passing time   & $8$  & \blank{2.6cm} \\
\end{tabularx}

\vspace{6pt}
\begin{center}
\begin{tikzpicture}[x=1cm, y=0.075cm]
  \foreach \y in {5,15,25,35}
    \draw[linegray, very thin] (0,\y) -- (10.4,\y);
  \foreach \y in {10,20,30,40}
    \draw[linegray] (0,\y) -- (10.4,\y);
  \foreach \y in {0,10,20,30,40}
    \node[left, font=\scriptsize] at (-0.05,\y) {\y\%};
  \foreach \x in {2.6,5.2,7.8}
    \draw[linegray, dashed] (\x,0) -- (\x,45);
  \foreach \i/\lab in {1/New library books, 2/Picnic tables, 3/Water bottle stations,
                       4/Longer passing time}
    \node[below, font=\scriptsize, align=center, text width=2.4cm]
      at ({2.6*\i-1.3},-0.5) {\lab};
  \draw[cerulean, thick] (0,45) -- (0,0) -- (10.4,0);
  \keyonly{\foreach \i/\h in {1/35, 2/25, 3/30, 4/10}
    \fill[keyred, opacity=0.55] ({2.6*\i-2.0},0) rectangle ({2.6*\i-0.6},\h);}
\end{tikzpicture}
\end{center}

\vspace{2pt}
The council's first draft item read: \emph{``Don't you agree the grant should buy
new books and picnic tables?''} Name the \textbf{two} instrument rules it breaks.

\par\writeline
\par\writeline

\end{enumerate}
\normalsize

% ======================================================================
\boxguard[10]
\parthead{Part D --- Extended Response \hfill 3 questions, 10 points each}

\small
\begin{enumerate}[leftmargin=2.2em, itemsep=10pt, topsep=2pt, start=31]

% ---------------------------------------------------------------- D1 (1.5-1.7)
\item Elmhurst Public Library has run \textbf{two} studies of its summer reading
program.

\vspace{1pt}
\textbf{Study A.} Library records were checked. Of the children whose families
\emph{chose} to sign up, $84\%$ met their reading goal; of those who did not sign
up, $52\%$ did. \\
\textbf{Study B.} Of $120$ volunteer children, a generator \emph{assigned} $60$ to
the program and $60$ to no program. Librarians collected the data by
\emph{observing} each child read aloud: $75\%$ of the program group met the goal
and $60\%$ of the others did.

\vspace{3pt}
(a) Which study is observational and which is an experiment? Give the one feature
that decides it --- and say why Study B's \emph{observing} does not change your answer.

\par\writeline
\par\writeline

(b) Study A's gap is $32$ points and Study B's is $15$. Which study lets the library
write \emph{``the program caused more children to meet their goal''}? Justify it.

\par\writeline
\par\writeline

(c) Explain why Study A's gap came out so much \emph{bigger} than Study B's.

\par\writeline
\par\writeline

% ---------------------------------------------------------------- D2 (1.4)
\boxguard[26]
\item \textbf{Brookside Apartments} mailed a survey about a proposed dog park to
all $720$ households. $180$ were returned, and $126$ of those said yes. The
management office's own sign-up records show that $288$ of the $720$ households
actually want one.

\vspace{2pt}
(a) Find the response rate, the survey's percent saying yes, and the number of
households that percent points to.

\begin{work}
  \dfrac{180}{720} &= 0.25 = 25\% \text{ response rate} \\
  \dfrac{126}{180} &= 0.70 = 70\% \\
  0.70 \times 720 &= 504 \text{ households}
\end{work}

(b) The records say the truth is $288$ of $720$, or $40\%$. Name the bias, and give
its size in percentage points \emph{and} in households.

\par\writeline
\par\writeline

(c) The manager mails again and gets $360$ returns, $252$ of them yes. Compute that
percent and explain what the result shows.

\begin{work}
  \dfrac{252}{360} &= 0.70 = 70\%
\end{work}

\par\writeline

(d) Describe one change that would actually reduce the bias, and name the bias it
removes.

\par\writeline
\par\writeline

% ---------------------------------------------------------------- D3 (1.3/1.7/1.8)
\boxguard[24]
\item \textbf{Granite Falls Middle School} ($840$ students) wants to know whether
to start a school garden club. The office has a numbered list of all $840$ students,
and the three grades are believed to differ a lot in their answers.

\vspace{2pt}
(a) Choose \textbf{one} of the four sampling techniques and justify it from what this
context allows.

\par\writeline
\par\writeline

(b) Name one of the five \textbf{data collection methods} and say why it fits this
question.

\par\writeline
\par\writeline

(c) Instead, a teacher asks her own advisory of $30$ students; $12$ say yes. A random
sample of the school later finds $38\%$. She reports: \emph{``My advisory got $40\%$
--- within $2$ points --- so asking one advisory works as well as a random sample.''}
Explain what is wrong with her claim, and state what she \emph{is} allowed to say.

\par\writeline
\par\writeline
\par\writeline

\end{enumerate}
\normalsize

\end{document}
